\documentclass[../../main.tex]{subfiles}
\graphicspath{{\subfix{../../image/}}}

\begin{document}

\section*{第四章小结}\label{section:ch04-小结}

本章主要讨论了随机变量（或分布函数）的数字特征与特征函数，它们都是概率分布的某种表征。这些讨论不但深化了对随机变量的认识，同时也为以后的研究作了必要的准备。

数字特征是描述随机变量特征的有效工具，它虽然不像分布函数那样完整地描述了随机变量，但是却具有很多优点：它较集中地反映了随机变量变化的一些平均特征（事实上数字特征多是某种平均值）；其次，大部分最重要的分布函数都由一、两个数字特征完全确定，而数字特征较易求得；最后，最常用也是最重要的几个数字特征——数学期望、方差、相关系数——都有明确的概率意义，同时又具有良好的性质。所有这些都决定了数字特征在概率论与数理统计中的重要地位。

特征函数与分布函数一一对应，它虽然不像分布函数那样有直观的概率意义，但却有好得多的分析性质，因此它是解决某些分布问题的有力工具，特别在处理独立随机变量和的分布问题上占有极重要的地位。关于特征函数与分布函数关系的另一些方面，将在下章中进一步讨论。

数字特征与特征函数都可以看作是分布函数的某种变换，而母函数则是概率论中引进的第一个变换，它在概率论发展史上有很大作用。到了现在，母函数以它的简明性以及在某些场合的有力应用而保持了它一定的地位。母函数与特征函数的性质与用途相似，因此本书采用互见的叙述方式以节省篇幅。

我们用特征函数作为工具研究了多元正态分布，建立了它的许多重要性质，这些良好性质从一个方面决定了正态分布在概率论、随机过程论与数理统计中的主角地位，正态分布重要性的另一方面——常见性——将在下章中得到深入的讨论。正态分布的大量特性只有在多元场合才能显现，因此可以说：只有懂得多元正态分布才算懂得正态分布。

最近几十年来，信息，信息的传输、变换与处理的概念正在渗入各种各样的学科，大大加深了人类对客观世界的认识。信息的概念与随机试验、概率分布等概率论概念密切相关，作为不肯定性度量的熵可以看作是概率空间的某种数字特征，因此我们对这些重要概念作了初步的介绍。

本章的作用是承上启下。本章引进的两个基本概念——数字特征与特征函数——将在极限定理的研究中起重要作用。

\end{document}